\documentclass[letterpaper]{article}
\usepackage[submission]{aaai2027}
\usepackage[hyphens]{url}
\usepackage{graphicx}
\urlstyle{rm}
\def\UrlFont{\rm}
\usepackage{natbib}
\usepackage{caption}
\usepackage{booktabs}
\usepackage{amsmath}
\frenchspacing
\pdfinfo{/TemplateVersion (2027.1)}
\setcounter{secnumdepth}{0}

\title{When Does Typed Conflict Control Help? Falsification-Guided Retrieval and Its Observed Boundaries}
\author{Anonymous Submission}
\affiliations{}

\begin{document}
\maketitle

\begin{abstract}
Retrieval-augmented language agents typically seek evidence that supports an emerging answer, while reflection remains weakly grounded in retrievable counter-evidence. We introduce Falsification-Augmented Retrieval (FAR), which decomposes an answer into dependent claims, assigns typed evidence requirements, retrieves support, refutation, and boundary evidence, and exposes a typed revision trace. On a 60-item Qwen3.5 9B machine-audited development diagnostic, typed control improves answer correctness over an otherwise matched untyped ablation by 0.078 (95\% paired bootstrap [0.034, 0.124]). This signal does not transfer into an end-to-end advantage on 350 upstream-labelled RAMDocs development items in either of two preregistered rounds. A frozen, zero-model-call attribution of the 226 shared errors finds 103 cases where conflict was detected but revision remained wrong; complete retrieval and partial-credit scoring do not restore a FAR advantage. All 34 positive per-item typed deltas on the Qwen diagnostic pass through changed revisions, even though removing typed revision improves the overall mean. Typed revision therefore mediates heterogeneous benefit and harm rather than a uniformly positive correction stage. We claim a narrow, machine-audited mechanism signal and characterize its observed boundary; we do not claim human-validated gold, external blindness, or general end-to-end superiority.
\end{abstract}

\section{Introduction}
Retrieval-augmented generation (RAG) grounds language models in external memory \cite{lewis2020rag}, but its default control loop is asymmetric: a query retrieves relevant passages, and generation accumulates passages that support an answer. Even iterative systems often ask whether more information is useful rather than what observation would falsify a particular claim. This asymmetry matters when the first answer has the wrong date, confuses a subsidiary with its parent, repeats an obsolete number, or upgrades correlation to causation.

We study \emph{falsification-guided retrieval}: given an initial answer, actively formulate retrievable evidence needs that could show where it is wrong. The central hypothesis is not merely that counter-evidence helps. It is that a \emph{typed evidence conflict} can serve as a control signal for both the next query and the revision action. A temporal conflict asks for version boundaries and repairs a timeline; a numerical conflict asks for alternative measurements and replaces or qualifies a value; a causal conflict seeks confounding or association-only evidence and downgrades the causal claim.

FAR makes this loop explicit. It (1) builds a claim dependency graph, (2) assigns positive typed evidence requirements, (3) generates support, refutation, and boundary queries, and (4) applies a typed revision with an auditable trace. We contribute an evaluation protocol, FalsiRAG-Bench, that places counter-evidence inside the retrieval corpus and scores whether it is found and used correctly. Our contributions are:
\begin{itemize}
    \item an explicit typed-conflict control framework with auditable claim, evidence, conflict, and revision traces, evaluated against a matched untyped control;
    \item a reproducible machine-audited synthetic benchmark and evaluation contract that keeps positive, negative, and mixed component evidence visible; and
    \item an applicability-boundary account that combines the positive Qwen development signal with two preregistered RAMDocs null results and frozen mechanism attribution.
\end{itemize}
The empirical claim is intentionally narrower than the full architecture: the present evidence supports typed conflict control, not a positive marginal contribution from every query family or revision component.

\section{Related Work}
Active retrieval and post-hoc repair are not new. FLARE retrieves from predicted upcoming sentences when generation confidence is low \cite{jiang2023flare}; Self-RAG learns reflection tokens for retrieval and critique \cite{asai2024selfrag}; and CRAG triggers corrective retrieval after assessing retrieval quality \cite{yan2024crag}. RARR researches attribution for generated claims and edits unsupported text while preserving the draft \cite{gao2023rarr}. Chain-of-Verification asks independent verification questions \cite{dhuliawala2024cove}, while RQ-RAG learns query rewriting, decomposition, and disambiguation \cite{chan2024rqrag}.

CounterRefine is the closest direct neighbor: it explicitly treats a draft answer as a hypothesis, performs answer-conditioned counterevidence retrieval, and applies a validated KEEP/REVISE repair step \cite{huang2026counterrefine}. FAR therefore makes no priority claim for counterevidence retrieval or answer revision alone. Its narrower hypothesis is that one explicit conflict ontology should jointly select three query families and claim-specific revision actions over a dependency graph. RAMDocs and MADAM-RAG instead study ambiguity, misinformation, and noise in already retrieved conflicting documents through multi-agent debate \cite{wang2025conflicting}. FalsiRAG-Bench isolates whether typed controls retrieve and use planted counter-evidence. FEVER supplies separately licensed external annotation candidates \cite{thorne2018fever}; machine seeds are never promoted to gold.

\section{Falsification-Augmented Retrieval}
Let an initial answer $A_0$ be decomposed into claims $C=\{c_i\}_{i=1}^{n}$ and dependency edges $D$. FAR produces a revised answer $A_1$, evidence map $M$, typed conflicts $K$, and revision trace $T$.

\subsection{Claim Graph}
Each node stores text, claim type, entities, numbers, time expressions, confidence, and verifiability. An edge $(c_j,c_i)\in D$ means that $c_i$ depends on $c_j$. The implementation rejects missing dependencies and cycles and processes nodes in topological order. Structured model decomposition is accepted only if it forms a valid graph and preserves at least 80 percent of source content tokens; otherwise a conservative deterministic decomposition is used.

\subsection{Typed Evidence Requirements}
For claim $c_i$, an assigner emits requirements $R_i$. The primary requirement follows claim semantics (temporal, entity, numerical, causal, or definition), while source reliability and counter-evidence are always considered. This is a forward retrieval contract, not a label attached after retrieval.

\subsection{Three-Way Counterfactual Queries}
For each primary type $t_i$, FAR generates
\begin{equation}
Q_i=\{q_i^{+},q_i^{-},q_i^{\partial}\},
\end{equation}
where $q_i^{+}$ seeks corroboration, $q_i^{-}$ seeks an observation incompatible with the claim, and $q_i^{\partial}$ seeks scope, time, unit, definition, or comparability boundaries. Tactics depend on $t_i$: numerical refutation searches for revised values and methodologies; causal refutation searches for association-only findings, confounders, and alternative causes. Model-generated queries must contain every enabled family; invalid output falls back to deterministic typed templates.

\subsection{Typed Conflict and Revision}
Retrieved passages are compared with the claim using high-precision rules or the layered VeraRAG conflict graph. FAR maps detected relations to seven control types. A fixed priority policy chooses an action: correct time, requalify entity, replace number, downgrade causality, prefer reliable source, clarify definition, retract, or qualify uncertainty. An optional language model realizes the selected action but cannot select a different policy. Every trace records before and after text, action, conflict types, evidence IDs, rationale, and confidence. No retrieved counter-evidence is treated as absence of a detected conflict, not proof of the claim.

\section{FalsiRAG-Bench}
The benchmark schema contains question, erroneous initial answer, atomic claims and dependencies, construction-reference evidence, counter-evidence with the refuted claim, typed conflict, expected revision, split, provenance, and annotation status.

\begin{table}[t]
\centering
\small
\begin{tabular}{lr}
\toprule
Candidate category & Count \\
\midrule
Temporal shift & 60 \\
Numerical conflict & 60 \\
Entity confusion & 60 \\
Causal overclaim & 60 \\
Multi-source conflict & 60 \\
\bottomrule
\end{tabular}
\caption{The deterministic v0.2 candidate build. Counts are construction facts, not human-validated gold statistics.}
\label{tab:data}
\end{table}

The current build derives evidence-backed candidates from VeraBench v1.1.2, with controlled perturbations and explicitly marked synthetic causal-boundary or low-reliability passages. It contains 175 documents, including 120 controlled synthetic documents. A source-document dependency group is assigned wholly to train, development, or test; the realization is 182/60/58 and has zero group leakage. The present study reports only the 60-item development split. A separate five-field test bundle is frozen for future work but is not treated as externally blind evidence.

All 300 labels are \texttt{machine\_seeded}; they are not human-validated gold. We audit them with two distinct signals that are blind to construction roles: Qwen2.5 7B structured preannotation (299/300 non-fallback outputs) and deterministic weak supervision (211/300 non-abstentions). Both signals exactly confirm the construction conflict type and action for 178 items; 122 items are disclosed as machine-disputed rather than silently relabeled. This machine agreement is quality evidence, not human inter-annotator agreement or adjudication. The construction check obtains 0.91 lexical counter-evidence recall@10 when pooling all three query families; this verifies corpus retrievability, not FAR performance.

\section{Experimental Design}
We evaluate a locally hashed Qwen3.5 9B model with matched retrieval settings. Baselines are Vanilla RAG, Multi-query RAG, Reflective RAG, closed-corpus CRAG-style and CounterRefine-style reproductions, and a Self-RAG-style inference approximation. The three style names explicitly avoid claiming equivalence to official trained or web-retrieval systems. This is a single-model development diagnostic; it does not establish multi-model generality.

Ablations remove typed conflict (untyped contradiction and generic queries), refutation queries, boundary queries, or typed revision. The primary gate compares typed and untyped FAR on aligned examples. Metrics are answer soft F1, unsupported claim rate, evidence precision, typed conflict F1, revision accuracy, overclaim reduction, and counter-evidence recall. We report category-stratified bootstrap confidence intervals, dependency-group sensitivity intervals, paired bootstrap differences, and exact McNemar tests for revision success.

Power is treated as a gate on claim strength rather than repaired after observing results. A preregistered simulation using the Qwen discordance pattern estimates only 0.414 exact-McNemar power for the planned three-family, 60-pair-per-family reproduction at a +0.078 effect. That study is therefore designated a directional reproduction: a nonsignificant result cannot establish absence, and family-level directions must be disclosed. The same retrospective audit shows that the RAMDocs gate had a confidence-interval half-width near three percentage points, too coarse to adjudicate smaller effects.

\begin{table}[t]
\centering
\small
\resizebox{\columnwidth}{!}{\begin{tabular}{lrrrr}
\toprule
Method & Answer & Typed F1 & Revision & Counter Recall \\
\midrule
far & 0.797 & 0.420 & 0.217 & 0.983 \\
vanilla & 0.725 & 0.000 & 0.000 & 0.667 \\
multi\_query\_rag & 0.701 & 0.000 & 0.000 & 0.817 \\
reflective\_rag & 0.678 & 0.000 & 0.000 & 0.750 \\
crag\_style\_reproduction & 0.772 & 0.000 & 0.000 & 0.667 \\
self\_rag\_style\_reproduction & 0.751 & 0.000 & 0.000 & 0.817 \\
counterrefine\_style\_reproduction & 0.710 & 0.000 & 0.000 & 0.883 \\
\bottomrule
\end{tabular}
}
\caption{Machine-seeded Qwen development diagnostic ($n=60$). ``Style'' baselines are scoped reproductions. These are not blind-test or human-gold results.}
\label{tab:main}
\end{table}

\begin{table}[t]
\centering
\small
\resizebox{\columnwidth}{!}{\begin{tabular}{lrrrr}
\toprule
Method & Answer & Typed F1 & Revision & Counter Recall \\
\midrule
far & 0.797 & 0.420 & 0.217 & 0.983 \\
minus\_typed\_conflict & 0.719 & 0.000 & 0.000 & 0.983 \\
minus\_refutation\_query & 0.826 & 0.429 & 0.267 & 0.967 \\
minus\_boundary\_query & 0.801 & 0.418 & 0.233 & 0.983 \\
minus\_typed\_revision & 0.873 & 0.423 & 0.000 & 0.983 \\
\bottomrule
\end{tabular}
}
\caption{Component ablations on the same diagnostic. The mixed results motivate our narrowed claim.}
\label{tab:ablations}
\end{table}

\section{Diagnostic Results}
FAR reaches 0.797 answer correctness, above the strongest baseline value of 0.772 from the CRAG-style reproduction, and retrieves counter-evidence for 0.983 of examples (Table~\ref{tab:main}). Its absolute revision accuracy is only 0.217, so the result is not evidence of a generally reliable automatic correction system.

The strongest controlled result is typed versus untyped conflict. FAR improves answer correctness by 0.078; reversing the sign reported by the ablation-centered evaluator gives a 95\% paired stratified bootstrap interval of [0.034, 0.124]. Typed-conflict F1 changes from 0 to 0.420 and revision accuracy from 0 to 0.217. Exact McNemar tests on revision success and typed-conflict correctness give $p=0.00024$ and $p<10^{-9}$, respectively. This supports typed conflict as a useful control signal on this diagnostic.

As a post-hoc label-audit sensitivity check, we stratify the same paired comparison by machine-consensus disposition. The answer-correctness difference remains positive on the machine-confirmed subset ($n=35$), +0.101 with 95\% bootstrap [0.039, 0.161], and on the machine-disputed subset ($n=25$), +0.047 [0.001, 0.094]. Thus the aggregate direction is not produced only by disputed labels. The subsets are small and category-imbalanced, and their dispositions are machine signals against the same construction reference; this is a sensitivity analysis, not independent label validation.

The remaining ablations contradict a stronger component-wise story. Removing refutation queries does not reduce answer correctness (0.826 versus 0.797), despite a small counter-evidence recall decrease. Removing boundary queries is also neutral on answer correctness (0.801 versus 0.797). Removing typed revision increases answer correctness to 0.873 while reducing revision accuracy and revision-action correctness to zero. Typed revision therefore provides explicit revision behavior at an answer-score cost; it is not supported as an unconditional accuracy improvement. We report these outcomes as negative or mixed evidence rather than attributing gains to all three query families or every repair stage.

The registered post-stop-rule attribution sharpens this trade-off. Among the
226 RAMDocs items missed by both FAR and Multi-Query RAG, the earliest-failure
buckets are 4 retrieval misses, 72 undetected conflicts, 103 detected conflicts
followed by wrong revision, 35 incomplete answer sets, and 12 overfull answer
sets; the format-mismatch bucket is empty. On the 179 items with complete
correct-document retrieval, FAR remains below the baseline by 0.011 (95\% CI
$[-0.050,0.028]$). Descriptive collection F1 also changes by $-0.003$ (95\% CI
$[-0.020,0.013]$), so strict exact match is not hiding a partial-set gain.
Although conflict detection is lower on RAMDocs (0.449) than on the Qwen
diagnostic (0.883), the detected subset has only a $+0.006$ paired difference
with a confidence interval crossing zero. Accordingly, all four preregistered
mechanism hypotheses are not supported.

The component path analysis also rejects our detection-only prediction. FAR has
34 positive continuous answer-score deltas over the untyped ablation, and every
one occurs on a changed-revision path; none occurs after detection without a
changed revision. Together with the higher mean obtained when typed revision is
removed, this identifies revision as a heterogeneous mediator of both local
benefit and larger aggregate harm. It does not identify a uniformly beneficial
component or rescue the RAMDocs gate.

Table~\ref{tab:preconditions} turns these negative and mixed findings into an
operational applicability checklist. These are observed necessary checks, not
sufficient conditions or a causal theorem: satisfying one row did not by itself
restore a FAR advantage.

\begin{table}[t]
\centering
\small
\begin{tabular}{p{0.25\columnwidth}p{0.65\columnwidth}}
\toprule
Check before use & Evidence and implication \\
\midrule
Correct evidence is retrievable & Complete correct-document retrieval still gives a $-0.011$ FAR difference; reachability is necessary but insufficient. \\
Conflict is typed and detected & 72 shared errors are undetected, while the detected subset is only $+0.006$ with an interval crossing zero; detection creates an opportunity, not a guarantee. \\
Revision preserves the signal & 103 shared errors are detected then revised incorrectly. Every positive Qwen typed delta passes through a changed revision, yet removing typed revision raises the mean; revision is both mediator and risk. \\
Conclusions survive metric choice & Collection F1 changes by $-0.003$ with an interval crossing zero; partial credit must be disclosed but cannot replace the frozen gate. \\
\bottomrule
\end{tabular}
\caption{Observed applicability checks for typed conflict control. They delimit where the mechanism can be meaningfully attempted; none is sufficient for an end-to-end gain.}
\label{tab:preconditions}
\end{table}

On a visible 100-pair FEVER SUPPORTS/REFUTES transfer diagnostic, the heuristic and NLI detectors both obtain 0.72 accuracy. NLI raises conflict recall from 0.30 to 0.40, but the paired accuracy difference is 0 with a 95\% bootstrap interval of $[-0.05,0.05]$ and McNemar $p=1.0$. This does not validate typed FAR behavior outside the controlled benchmark.

\section{Mechanism Interpretation and Applicability Boundary}

\subsection{An opportunity chain, not a monotone pipeline}

Typed conflict control can influence an answer only through a sequence of
opportunities: relevant counter-evidence must be retrievable; the observed
relation must map to a usable conflict type; the detector must expose that
type; the policy must select an appropriate action; and the realized revision
must preserve correct content while changing the error. The completed evidence
does not license a multiplicative causal decomposition of these stages, but it
does show why success at one stage is not sufficient. Complete correct-document
retrieval does not restore a RAMDocs advantage, the detected-conflict subset is
not reliably positive, and detected-then-wrong revision is the largest shared
failure bucket. Typed control is therefore best understood as an opportunity
to intervene, not as a monotone improvement applied after retrieval.

\subsection{Why the local gain and external null can coexist}

FalsiRAG-Bench makes counter-evidence lexically reachable and constructs labels
in the same temporal, numerical, entity, causal, definition, and source-quality
vocabulary used by the controller. That alignment creates a diagnostic in
which the typed mechanism can be observed; it does not establish that the
mechanism transports. RAMDocs instead evaluates end-to-end answer-set recovery
under ambiguity, misinformation, and noise. Its two failed gates show that the
local signal is insufficient on that distribution. They do not, by themselves,
prove that construct alignment is the sole cause: the registered
H-conflict-shape test is also not supported. The defensible synthesis is thus
conditional rather than contradictory---the typed mechanism has a local signal,
while its transport conditions remain only partially identified.

\subsection{Revision is both carrier and risk}

Every observed positive Qwen typed-minus-untyped delta passes through a changed
revision, yet removing typed revision raises the aggregate answer score. This
combination rules out both a detection-only account and an always-revise policy.
It suggests a future selective-intervention design: revise only when the
expected benefit of the typed action exceeds its risk, and otherwise preserve
or abstain. No such selector has been evaluated here, so this is a design
hypothesis rather than a reported FAR capability. A proper test would register
the selector before execution and score selective risk, coverage, and answer
quality without replacing the current frozen gates.

\subsection{What the remaining workstreams can identify}

WS2 asks whether the local signal is specific to Qwen or directionally recurs
across Mistral, Gemma, and Llama under a fixed two-arm protocol. WS3 asks whether
typed effects track preregistered external conflict conditions rather than only
the construction-derived diagnostic. Until their independently verified
releases exist, neither cross-family recurrence nor an external applicability
condition is a result. Their role is to choose among the registered A/B/C paper
lines, not to rescue RAMDocs G-A, unlock held-out data, or reinstate a general
end-to-end claim.

\section{Limitations and Responsible Use}
The candidate set reuses a small controlled source corpus, includes synthetic passages, and has correlated variants. Its labels are not human-validated gold, the reported split is development data, and evaluation is not externally blind. The machine audit measures agreement with construction labels and cannot reveal errors shared by construction, the LLM signal, and deterministic rules. Current completed results come from one local model and do not yet establish multi-model generality. Rule conflict detectors are high precision but incomplete; model detectors inherit calibration and domain-shift risk. Retrieval failure can hide a true counterexample, while spurious low-quality contradictions can induce over-correction. The benchmark is not current factual advice and is unsuitable for medical, legal, financial, or policy decisions. Final authors remain responsible for checking all text, citations, experiments, and disclosure requirements.

Our preregistered no-budget 2+4 profile completed only its external Phase A and
then stopped. Round 1 on 350 RAMDocs development items with upstream labels gave
FAR and the strongest baseline, Multi-Query RAG, the same 0.311 strict exact
match. Round 2 changed only FAR's final answer consolidation layer; its G-A
comparison also failed, with FAR at 0.309 versus 0.311 for the frozen
Multi-Query baseline, paired difference $-0.0029$, bootstrap 95\% CI
$[-0.0314,0.0286]$, and McNemar $p=1.0$. This second failed G-A triggers the
registered stop rule and downgrades the 2+4 branch to an
applicability-boundary analysis for typed conflict control. Phase B not run:
the DeepSeek/GLM/Meta jury, author-blind adjudication, G-K/G-S, jury rescoring,
and three-family model matrix were not run. Held-out not run: neither the
FalsiRAG-Bench held-out split nor the RAMDocs held-out split was evaluated.
The RAMDocs evidence uses upstream labels and is not human inter-annotator
agreement, not publication-grade human gold, and not an externally held test.
Any future attempt to revive this profile would have to start from the recorded
failure analysis, disclose the negative refutation and boundary ablations, the
typed revision trade-off, and the FEVER null result, and demonstrate transfer
across 7--9B open-model families before replacing this single-model diagnostic
claim.

\section{Conclusion}
FAR turns ``what could make this answer wrong?'' into a typed, retrievable control protocol. On the released machine-audited Qwen development diagnostic, typed conflict control outperforms a matched untyped alternative, but two preregistered RAMDocs rounds and the frozen attribution reject a general end-to-end or upstream-only explanation. Revision is the path through which all observed positive typed deltas occur and also the component whose removal raises the aggregate score. The resulting contribution is a narrow, reproducible account of revision-mediated heterogeneous effects and their observed boundary---not human-gold, blind-test, or general superiority evidence.

\bibliography{references}
\end{document}
